\abstract{РЕФЕРАТ}

Объем работы равен \formbytotal{lastpage}{страниц}{е}{ам}{ам}. Работа содержит \formbytotal{figurecnt}{иллюстраци}{ю}{и}{й}, \formbytotal{tablecnt}{таблиц}{у}{ы}{} 
и \arabic{bibcount} библиографических источников. Количество приложений – 2. Графический материал представлен в приложении А. Фрагменты исходного кода представлены в приложении Б.

Перечень ключевых слов: Python, REST API, FastAPI, PostgreSQL,
информационно-аналитическая система, проверка учебно-методических материалов, веб-приложение, генерация титульных листов, PDF-отчёты.

Объектом разработки является информационно-аналитическая система автоматизации проверки и обработки учебно-методических материалов редакционно-издательского отдела.

Целью работы является разработка информационно-аналитической системы для автоматизации процессов проверки, оформления и сопровождения учебно-методических документов.

Разработанная система реализована на языке Python с использованием web-фреймворка FastAPI и СУБД PostgreSQL. Клиентская часть разработана с применением HTML, CSS и JavaScript. Взаимодействие между компонентами системы осуществляется через REST API.

Архитектура системы построена по клиент-серверной модели. Серверная часть выполняет обработку запросов пользователей, проверку DOCX-документов, генерацию титульных листов, формирование PDF-отчётов и взаимодействие с базой данных. Клиентская часть предоставляет web-интерфейс для загрузки документов и просмотра результатов проверки.

В системе реализован модуль автоматической проверки DOCX-документов, выполняющий анализ параметров форматирования и структуры учебно-методических материалов. Результаты проверки используются для формирования PDF-отчётов с перечнем найденных ошибок..

Разработанная система позволяет сократить время проверки документов, уменьшить количество ошибок оформления и повысить эффективность работы редакционно-издательского отдела.



\selectlanguage{english}
\abstract{ABSTRACT}

The volume of work is \formbytotal{lastpage}{page}{}{s}{s}. The work contains \formbytotal{figurecnt}{illustration}{}{s}{s},\formbytotal{tablecnt}{tables}{}{s}{s} and \arabic{bibcount} bibliographic sources. The number of applications is 2. The graphic material is presented in Appendix A. Fragments of the source code are presented in Appendix B.

List of keywords: Python, REST API, FastAPI, PostgreSQL,
information and analytical system, verification of teaching materials, web application, generation of cover pages, PDF reports.

The object of the development is an information and analytical automation system for checking and processing educational and methodological materials of the editorial and publishing department.

The purpose of the work is to develop an information and analytical system for automating the processes of verification, registration and maintenance of educational and methodological documents.

The developed system is implemented in Python using the FastAPI web framework and the PostgreSQL database management system. The client side is developed using HTML, CSS, and JavaScript. The interaction between the system components is carried out through the REST API.

The architecture of the system is based on the client-server model. The server side processes user requests, validates DOCX documents, generates title pages, generates PDF reports, and interacts with the database. The client side provides a web interface for downloading documents and viewing the verification results.

The system implements an automatic verification module for DOCX documents, which analyzes the formatting parameters and structure of teaching materials. The verification results are used to generate PDF reports with a list of errors found..

The developed system makes it possible to reduce the time needed to check documents, reduce the number of design errors, and improve the efficiency of the editorial and publishing department.
\selectlanguage{russian}